\appendix

\chapter{原证明、计算证据与本篇的对应}

本附录以稳定来源编号连接讲义与原研究。压缩包的 sources 目录保持原项目相对路径；来源散列见 SOURCE\_MANIFEST.json。正文按教学顺序重写，原始计数表、推导及历史实验仍保留在对应原文中。

技术原稿中的“未解”或“下一步”可能属于历史阶段；本篇第十四讲按后续定理更新其覆盖范围。核验记录列出原核验的证据等级，本轮没有把所有历史程序重新运行一遍。

原验证器的命令应在原项目根目录执行。证据快照便于检查与追溯，少数历史程序仍可能依赖原工作区；本篇PDF和LaTeX编译不依赖它们。

\section{S01：合数平均的临界尺度：最早成立、最终阈值与无条件网络}\label{source:S01}

M/H/N、赋值下界、B/C猜想、四步网络、八平均十三元；历史状态由S02更新。对应本篇第1—3、6、14讲。

\par\smallskip{\small 原路径：\nolinkurl{outputs/general_arity/composite_critical_scale_and_conjecture.md}}\par

核验登记：{\small\nolinkurl{ten-twelve-seventeen-complete}}。原证据等级：{\small\nolinkurl{theorem-certificate}}。

核验登记：{\small\nolinkurl{seven-average-fifteen-complete}}。原证据等级：{\small\nolinkurl{theorem-certificate}}。

核验登记：{\small\nolinkurl{six-average-ten-complete}}。原证据等级：{\small\nolinkurl{theorem-certificate}}。

核验登记：{\small\nolinkurl{composite-critical-scale}}。原证据等级：{\small\nolinkurl{symbolic-check}}。

核验登记：{\small\nolinkurl{eight-average-thirteen-complete}}。原证据等级：{\small\nolinkurl{theorem-certificate}}。

\section{S02：GRH主线三项任务：本轮完成范围与当前剩余}\label{source:S02}

最近一般元数汇合范围；三项尚未完成的任务。对应本篇第1、10—14讲。

\par\smallskip{\small 原路径：\nolinkurl{outputs/general_arity/general_k_three_tasks_progress_20260915.md}}\par

本文件没有单独匹配的核验登记；其书面推导或依赖结果须在原文中核对，不能据此虚构独立程序。

\section{S03：接受GRH后的完整平方族定理与剩余证明任务}\label{source:S03}

条件性平方族、无条件子族、GRH接口与未闭合量词。对应本篇第8—11、14讲。

\par\smallskip{\small 原路径：\nolinkurl{outputs/general_arity/general_k_grh_stage_and_remaining_tasks_20260915.md}}\par

本文件没有单独匹配的核验登记；其书面推导或依赖结果须在原文中核对，不能据此虚构独立程序。

\section{S04：任意平均元数的线性阈值：非单位差分也是可执行资源}\label{source:S04}

非单位差分、中点碰撞、一般EGZ、统一尾部与4/6造零收尾。对应本篇第5讲。

\par\smallskip{\small 原路径：\nolinkurl{outputs/general_arity/arbitrary_arity_linear_threshold.md}}\par

核验登记：{\small\nolinkurl{arbitrary-arity-linear-threshold}}。原证据等级：{\small\nolinkurl{symbolic-check}}。

\section{S05：合数平均：四平均、六平均与原位置副本转译}\label{source:S05}

秩障碍、副本与半块转译、张量网络、全有理输入引理、摘要误用撤回。对应本篇第3—4、13讲。

\par\smallskip{\small 原路径：\nolinkurl{outputs/general_arity/composite_arity_four_six_transfer.md}}\par

核验登记：{\small\nolinkurl{six-average-nine-fixed-network}}。原证据等级：{\small\nolinkurl{theorem-certificate}}。

核验登记：{\small\nolinkurl{composite-four-six-transfer}}。原证据等级：{\small\nolinkurl{symbolic-check}}。

\section{S06：从素数完整证明到合数临界问题：策略、可迁移定理与真实困难}\label{source:S06}

素数因子块均值转译、全部三倍以上维数及合数特有困难。对应本篇第4、14讲。

\par\smallskip{\small 原路径：\nolinkurl{outputs/general_arity/averaging_proof_strategy_and_composite_frontier_20260915.md}}\par

核验登记：{\small\nolinkurl{square-arity-controller-nine-thirteen}}。原证据等级：{\small\nolinkurl{symbolic-check}}。

核验登记：{\small\nolinkurl{prime-factor-block-mean-transfer}}。原证据等级：{\small\nolinkurl{symbolic-check}}。

\section{S07：四平均七元完整判据：\textbackslash{}((4,2,1)\textbackslash{}) 核与模7控制器}\label{source:S07}

4+2+1真实周期和M(4)=H(4)=7。对应本篇第6—7讲。

\par\smallskip{\small 原路径：\nolinkurl{outputs/general_arity/examples/four_average_seven_complete.md}}\par

核验登记：{\small\nolinkurl{four-average-seven-complete}}。原证据等级：{\small\nolinkurl{theorem-certificate}}。

\section{S08：四平均九元完整判据：一个四步周期与模9控制器}\label{source:S08}

偶元数端点种子与模9完整陪集；后由S18推广。对应本篇第6、12讲。

\par\smallskip{\small 原路径：\nolinkurl{outputs/general_arity/examples/four_average_nine_complete.md}}\par

核验登记：{\small\nolinkurl{four-average-nine-complete}}。原证据等级：{\small\nolinkurl{theorem-certificate}}。

\section{S09：六平均九元：所有输入固定四步完成}\label{source:S09}

固定四步网络及固定词三步下界。对应本篇第3讲。

\par\smallskip{\small 原路径：\nolinkurl{outputs/general_arity/examples/six_average_nine_fixed_network.md}}\par

核验登记：{\small\nolinkurl{six-average-nine-fixed-network}}。原证据等级：{\small\nolinkurl{theorem-certificate}}。

\section{S10：六平均十元完整G判据：第一个非自动的算术成功维数}\label{source:S10}

G=1或2、模5控制器、两步返回与H(6)=10。对应本篇第1、6讲。

\par\smallskip{\small 原路径：\nolinkurl{outputs/general_arity/examples/six_average_ten_complete.md}}\par

核验登记：{\small\nolinkurl{six-average-ten-complete}}。原证据等级：{\small\nolinkurl{theorem-certificate}}。

\section{S11：十平均与十二平均17元完整G判据，以及原B公式中的自动维数}\label{source:S11}

两套实际返回、模17尺度、自动B点表、H(10)=H(12)=17。对应本篇第3、6、14讲。

\par\smallskip{\small 原路径：\nolinkurl{outputs/general_arity/ten_twelve_average_seventeen_complete.md}}\par

核验登记：{\small\nolinkurl{ten-twelve-seventeen-complete}}。原证据等级：{\small\nolinkurl{theorem-certificate}}。

\section{S12：十平均十六元与一族确定性三块网络}\label{source:S12}

固定2b+2步网络族、十平均十六元。对应本篇第3讲。

\par\smallskip{\small 原路径：\nolinkurl{outputs/general_arity/ten_average_sixteen_and_dyadic_network_family.md}}\par

核验登记：{\small\nolinkurl{ten-sixteen-dyadic-network-family}}。原证据等级：{\small\nolinkurl{theorem-certificate}}。

\section{S13：平方元数的统一正向控制器与九平均十三元}\label{source:S13}

平方核心全参数周期、两步数字返回、除二与完整根。对应本篇第7讲。

\par\smallskip{\small 原路径：\nolinkurl{outputs/general_arity/square_arity_uniform_controller_and_nine_thirteen.md}}\par

核验登记：{\small\nolinkurl{square-arity-controller-nine-thirteen}}。原证据等级：{\small\nolinkurl{symbolic-check}}。

\section{S14：一般k平均规划评议：平方族入口可直接补齐，尺度任务可以数论化}\label{source:S14}

全t平方入口、H\_t、一步终端尺度限制。对应本篇第8、13讲。

\par\smallskip{\small 原路径：\nolinkurl{outputs/general_arity/general_k_plan_review_and_square_entry_20260915.md}}\par

核验登记：{\small\nolinkurl{square-core-entry-review}}。原证据等级：{\small\nolinkurl{symbolic-check}}。

\section{S15：对一般k规划建议的复核与采纳}\label{source:S15}

Gamma1扩环接口、正确一般整数核心与G公式。对应本篇第8、12讲。

\par\smallskip{\small 原路径：\nolinkurl{outputs/general_arity/general_k_review_response_20260915.md}}\par

核验登记：{\small\nolinkurl{square-core-entry-review}}。原证据等级：{\small\nolinkurl{symbolic-check}}。

\section{S16：H\_t文献调查：已有小生成元理论、无条件充分大结论与全范围边界}\label{source:S16}

小生成元、特征、Pollack、Bach/GRH、998参数证书、105及立方因子反例。对应本篇第9讲。

\par\smallskip{\small 原路径：\nolinkurl{outputs/general_arity/ht_literature_and_eventual_square_completion_20260915.md}}\par

核验登记：{\small\nolinkurl{ht-literature-interfaces}}。原证据等级：{\small\nolinkurl{theorem-certificate}}。

\section{S17：统一四／六阶控制器完成2、3奇次幂临界点}\label{source:S17}

2s²/3s²完整控制器、508组小参数与条件性临界族。对应本篇第10讲。

\par\smallskip{\small 原路径：\nolinkurl{outputs/general_arity/two_three_odd_power_critical_grh_completion.md}}\par

核验登记：{\small\nolinkurl{two-three-odd-power-critical}}。原证据等级：{\small\nolinkurl{theorem-certificate}}。

\section{S18：全部偶数元数的2k+1端点：旧种子与统一取逆}\label{source:S18}

全部偶数元数2k+1端点；接奇数端点得到所有k。对应本篇第12讲。

\par\smallskip{\small 原路径：\nolinkurl{outputs/general_arity/even_arity_all_endpoints_completion.md}}\par

核验登记：{\small\nolinkurl{even-arity-all-endpoints}}。原证据等级：{\small\nolinkurl{symbolic-check}}。

\section{S19：奇合数元数的奇余量区间：按真实互素条件迁移旧证明}\label{source:S19}

奇k奇r互素带、奇素数幂全部奇中间维数。对应本篇第12讲。

\par\smallskip{\small 原路径：\nolinkurl{outputs/general_arity/odd_composite_coprime_middle_band_completion.md}}\par

核验登记：{\small\nolinkurl{odd-composite-coprime-middle-band}}。原证据等级：{\small\nolinkurl{symbolic-check}}。

\section{S20：较大素数高奇次幂：两种统一三步正向反射}\label{source:S20}

p>=5高奇次幂的两种三步反射，未得完整根；立方未覆盖。对应本篇第11讲。

\par\smallskip{\small 原路径：\nolinkurl{outputs/general_arity/large_prime_high_odd_power_three_step_reflections.md}}\par

核验登记：{\small\nolinkurl{large-prime-high-odd-power-reflections}}。原证据等级：{\small\nolinkurl{symbolic-check}}。

\section{S21：一般k平均的证明计划：先完成核心，再接成全部维数}\label{source:S21}

theta周期分类、临界／混合／接续三项任务及验收边界。对应本篇第10—14讲。

\par\smallskip{\small 原路径：\nolinkurl{outputs/general_arity/general_k_averaging_proof_plan_20260915.md}}\par

核验登记：{\small\nolinkurl{square-core-entry-review}}。原证据等级：{\small\nolinkurl{symbolic-check}}。

\section{S22：逆零和定理的原位置转译：碰撞修复与 2p+4 元七步族}\label{source:S22}

通用双碰撞修复及奇数元数七步二参数族；素数加权结果仅引用前史。对应本篇第13讲。

\par\smallskip{\small 原路径：\nolinkurl{outputs/general_arity/inverse_zero_sum_averaging_progress.md}}\par

核验登记：{\small\nolinkurl{inverse-zero-sum-transfer}}。原证据等级：{\small\nolinkurl{symbolic-check}}。

\section{S23：从十五元推广到合数端点：统一CRT证书与两素数完整入口}\label{source:S23}

36个CRT陪集证书的历史阶段及后续替代范围，不重教素数篇陪集算法。对应本篇第12、13讲。

\par\smallskip{\small 原路径：\nolinkurl{outputs/general_arity/composite_endpoint_crt_transfer.md}}\par

核验登记：{\small\nolinkurl{composite-endpoint-transfer}}。原证据等级：{\small\nolinkurl{theorem-certificate}}。

\section{S24：q-平均的根格、\textbackslash{}(S\textbackslash{})-算术与控制器猜想}\label{source:S24}

原子秩与压缩参数秩的区别；几何／正运输仅作为有边界的候选框架。对应本篇第13讲。

\par\smallskip{\small 原路径：\nolinkurl{outputs/general_arity/triple_average_q_average_arithmetic_geometry.md}}\par

本文件没有单独匹配的核验登记；其书面推导或依赖结果须在原文中核对，不能据此虚构独立程序。

\section{S25：从二平均到任意平均元数：当前全貌与临界猜想}\label{source:S25}

历史全局表与最新定理对照；不沿用过时未解项。对应本篇第14讲。

\par\smallskip{\small 原路径：\nolinkurl{outputs/general_arity/averaging_global_status_20260912.md}}\par

本文件没有单独匹配的核验登记；其书面推导或依赖结果须在原文中核对，不能据此虚构独立程序。

\section{S26：移动见证再合并：全部元数、全部超临界维数的同一入口}\label{source:S26}

所有k、n>=2k+1的移动见证双块入口，未提供全范围终止。对应本篇第12—13讲。

\par\smallskip{\small 原路径：\nolinkurl{research/prime_arity_padic/move_witnesses_then_merge.md}}\par

本文件没有单独匹配的核验登记；其书面推导或依赖结果须在原文中核对，不能据此虚构独立程序。

\section{S27：用载体的 p 进格位移审查统一机制}\label{source:S27}

素数单载体全字母表刚性、秩下降交条件、不能凭群或运输代替原位置词。对应本篇第13讲。

\par\smallskip{\small 原路径：\nolinkurl{research/prime_arity_padic/carrier_balance_and_uniform_mechanism.md}}\par

本文件没有单独匹配的核验登记；其书面推导或依赖结果须在原文中核对，不能据此虚构独立程序。

\section{S28：任意维数的新统一结果：复制稳定化、双随机矩阵与消去猜想}\label{source:S28}

第2及8节：任意k复制稳定化、正双随机证书及原位置消去缺口。对应本篇第13讲。

\par\smallskip{\small 原路径：\nolinkurl{outputs/triple_average/history/triple_average_stabilization_theorem.md}}\par

核验登记：{\small\nolinkurl{stabilization-theorem}}。原证据等级：{\small\nolinkurl{finite-sanity}}。

\section{S29：任意维数的三幂分块压缩：把一般输入降到对数秩标准核}\label{source:S29}

同时好划分的独立组合依据。对应本篇第4讲。

\par\smallskip{\small 原路径：\nolinkurl{outputs/triple_average/history/triple_average_power_partition_reduction.md}}\par

核验登记：{\small\nolinkurl{power-partition-compression}}。原证据等级：{\small\nolinkurl{finite-sanity}}。

\section{S30：H\_t文献来源与访问等级}\label{source:S30}

外部文献版本、已读范围和未取得全文的边界。对应本篇第9、来源附录讲。

\par\smallskip{\small 原路径：\nolinkurl{work/ht_literature_20260915/source_index.md}}\par

本文件没有单独匹配的核验登记；其书面推导或依赖结果须在原文中核对，不能据此虚构独立程序。
